\section{Ridwan and Malik}

Christians, particularly Christian monks or other holy men, often play a considerable role in the traditions concerning the martyrdom of the Imam Hussayn, usually acting as a rebuke to put to shame nominal Muslims who betrayed the Imam Hussayn. Here is a sample.

\begin{quotex}
Hussayn's head, according to popular tradition, was sent to Damascus with a large company of bodyguards to make sure it got to Yazid. On the way, the caravan stopped for a night below a hermitage where a Christian monk lived, spending his life in solitary worship. As they sat down for dinner, a hand wrote on the wall with letters of blood: ``Would a community that killed Hussayn hope for the intercession of his grandfather (Muhammad) on the day of Reckoning?'' The monk looked down and saw the writing on the wall and the head of the Imam Hussayn surrounded by an aura of bright light. He offered Ibn Saad ten thousand dinars to keep the head for a night. He took the head and, with it pressed to his bosom, spent the night weeping.\footnote{Ayoub, op. cit., p. 133.}
\end{quotex}


The special place which St. Peter holds in the Catholic Tradition is well known. He is called ``Prince of the Apostles'' and the ``keeper of the keys of the Kingdom.'' The popes are considered the successors of St. Peter. In this respect the Shi'ite tradition is also very interesting. Mahmoud Ayoub says:

\begin{quotex}
Ridwan (the keeper of paradise) and Malik (the keeper of hell) will come to the Prophet (Muhammad) and deliver into his hands the keys of paradise and hell. The Prophet will give them (the keys) to `Ali, who will then permit whomsoever he wishes to enter paradise and whomsoever he wishes to enter the fire (of hell). In this mood of exultation, the tradition concludes, ``And hell shall be on that day more obedient to `Ali than a young servant would be to his master.'' The intercessory character of this tradition is obvious. It is, however, interesting to note the similarity in this tradition between `Ali, the viceregent and successor of Muhammad, and Simon (St.) Peter, the prince of the apostles and keeper of the keys of the Kingdom. In the Shi`i doctrine of the Imamate-succession, St. Peter (\emph{Shamun al-Safa}) is declared to be a prophet.\footnote{Ibid., p. 203.}
\end{quotex}

Thus, not only is the role of `Ali Ibn Abi Talib in the Shi`a tradition very similar to that of St. Peter in the Catholic tradition, but in the Shi`a tradition as well as the Catholic tradition St. Peter is accorded a position superior to that of the other apostles of Jesus.

It would be very wrong to say that the role of Fatima al-Zahra (the Radiant), daughter of the prophet Muhammad, wife of `Ali Ibn Abi Talib, the First Imam, and mother of the Imams Hassan and Hussayn, is identical to the role of the Virgin Mary, mother of Jesus Christ, in the Catholic and Orthodox Churches. The expression ``Mother of God'' (Greek: \emph{Theotokos,} Latin: \emph{Mater Dei}) is simply unthinkable in an Islamic context. Nevertheless, there are parallels between the role of the Virgin Mary in the Catholic and Eastern Orthodox Churches on the one hand and the role of Fatima Zahra in Shi`a Islam on the other...

Jacob of Serug, known as ``flute of the Holy Spirit and harp of the faithful church,'' was born in 451 AD and died in 519 AD. Though born at Curtem on the Euphrates, he became Bishop of Serug, also in Syria. He is best known for his Syriac homilies on the Virgin Mary, such as this extract from one of his homilies:

\begin{quotex}
\emph{The story of Mary stirs in me, to show itself in wonder; you, wisely, prepare your minds!}

\emph{The Holy Virgin calls me today to speak of her; let us purge our hearing for her luminous tale, lest it be dishonored.}

\emph{Blessed of women, by whom the curse of the land was eradicated, and the sentence henceforth has come to an end.}

\emph{Modest, chaste and filled with beauties of holiness, so that my mouth is inadequate to speak a word concerning her.}

\emph{Second Eve who generated life among mortals, and paid and rent asunder the bill of Eve her mother.}

\emph{Virgin who without marital union marvelously became a mother, a mother who remained without change in her virginity.}

\emph{The image of her beauty is more glorious and exalted than my composition; I do not dare let my mind depict the form of her image.}

\emph{It is easier to depict the sun with its light and its heat than to tell the story of Mary in its splendor.}

\emph{Perhaps the rays of the sphere can be captured in pigments, but the tale concerning her is not completely told by those who preach... }

\emph{She is virgin and mother and wife of a husband yet unmated; how may I speak if I say that she is incomprehensible?}

\emph{Then He sent a Watcher (angel) from the Heavenly legions, that he might bring the good tidings to the blessed one, most fair.}

\emph{Gabriel, the great chief of the hosts, descended; he went down to her as he had been sent from God.}

\emph{Because she alone was worthy of the great mystery which was rich in divine revelations.}

\emph{With prayers and in limpidity and in simplicity, Mary received that spiritual revelation.}

\emph{She being holy and standing in wonder in God's presence, her heart was poured forth with love in prayer before Him.}

\emph{She was in prayer, as also Daniel was in prayer, when this same Watcher of light descended to him.}

\emph{The Watcher had descended while Mary was standing in prayer; he gave her the greeting which was sent to her from the Most High... }

\emph{``Hail Mary, our Lord is with you,'' he was saying to her, ``you will conceive and bear a son in your virginity.''}

\emph{She said to him: ``How will this be as you say, since I am a virgin and there is no fruit of virgins?''... }

\emph{The Watcher told her that she would conceive the Son of God, but she did not accept it until she was well informed... }

\emph{She asked the Watcher: ``How will what you tell me take place?'' He began explaining the way of the Son and His descent within her:}

\emph{``The Holy Spirit will come to you with solemnity, and the Power of the Most High will overshadow you, O most blessed one.''}

\emph{Here all speech of the tongue is superfluous; one does not speak except with the wonder of faith... }

\emph{In this way, the Watcher announced to her that he had come from the house of the Father: ``The Spirit will come and then the Power of the Most High will descend.''}

\emph{Indeed, the Holy Spirit came to Mary, to let loose from her the former sentence of Eve and Adam.}

\emph{He sanctified her, purified her and made her blessed among women; He freed her from the curse of sufferings on account of Eve her mother.}

\emph{The Spirit freed her from that debt that she might be beyond transgression when He solemnly dwelt in her.}

\emph{He purified the Mother by the Holy Spirit while dwelling in her, that He might take from her a pure body without sin... }

\emph{Mary appeared to us as a sealed letter in which were hidden the mysteries of the Son and His depth.}

\emph{She gave her body as a clean sheet; the Word wrote His essence on it, corporeally.}

\emph{She was the letter, not because she was sealed after she was inscribed, but the Divinity sealed her and then wrote on her... .}

\emph{With her the Father sent us tidings full of good things, and through her, forgiveness to all condemned for their bonds of sin.}

\emph{By her, emancipation was sent to Adam who had been enslaved; he became an heir and came in among the sons, as he had been.}

\emph{Because of her, confusion was lifted from womanhood; the reproach of all women passed away from the nations.}

\emph{Because of her, the way to Eden, which had been blocked, was opened; the serpent fled and men passed along it to God.}

\emph{Because of her, the Cherub had removed his lance that he might no longer guard the Tree of Life which offered itself to those who ate it.}

\emph{She gave us a sweet fruit, full of life, that we might eat from it and live forever with God.

\flright{\textsc{Jacob of Serug}, \emph{On the Mother of God}, translated from the Syriac by Mary Hansbury, introduction by Sebastian Brock (Crestwood, New York, 1998), pp. 18-42.}
\end{quotex}


Perhaps the most widespread hymn in the Eastern Orthodox Church is the seventh century Byzantine hymn usually known by its Greek name \emph{Akathist} (or \emph{Akafist} in Church Slavonic). The hymn is very long and deals with various topics. Here is a strophe which deals with the Virgin Mary:

\begin{quotex}
Rejoice, through you joy rings out again.

Rejoice, through you sorrow is put to flight.

Rejoice, O resurrection of fallen Adam.

Rejoice, O redemption of the tears of Eve.

Rejoice, O sublime peak of human intellect.

Rejoice, O profound abyss even for Angel eyes.

Rejoice, for in you the King's throne was elevated.

Rejoice, for you bear the One Who sustains everything.

Rejoice, O star that goes before the Sun.

Rejoice, O womb of the Incarnate God.

Rejoice, for through you all creation is renewed.

Rejoice, for through you the Creator became a baby.

Rejoice, O Virgin and Brid
\end{quotex}


In the eleventh century the Clunic monks composed one of the most widely known of Catholic hymns, generally known by its Latin title ``Salve Regina''. I do not have the original Latin at hand, but here is a translation:

\begin{quotex}
Hail Holy Queen, Mother of mercy!

Hail, our life, our sweetness and our hope.

To thee do we cry,

Poor banished children of Eve.

To thee do we send up our sighs,

Mourning and weeping in this valley of tears.

Turn then, most gracious advocate,

Thine eyes of mercy toward us.

And after this our exile

Show unto us the blessed fruit of thy womb, Jesus.

O clement, o loving, o sweet Virgin Mary!
\end{quotex}


The expression ``Mother of God,'' or \emph{Theotokos}, is simply inconceivable and unthinkable in an Islamic context, as we have noted above. However, there are certain similarities between the role of the Virgin Mary in Catholicism, particularly in Spain and Ireland, and that of Fatima, the daughter of Muhammad, called \emph{al-Zahra}, ``the Radiant,'' mother of the martyred Imams Hassan and Hussayn, the second and third Holy Shi`a Imams. Fatima al-Zahra appears as the \emph{Mater Dolorosa} (Sorrowful Mother), weeping for her martyred sons.\footnote{Abu Muhammad Ordoni, \emph{Fatima the Glorious} (Qum, Iran, 1992); Nasser Makorem Chirazi, \emph{La Dame La Plus Prestigieus de Monde}: ``Fatima -- Zahra'' (Qum, Iran, undated); Henry Corbin, \emph{Face de Dieu, Face de l'Homme} (Paris, 1983), pp. 149, 154-155; Moojan Momen, \emph{An Introduction to Shi`i Islam} (New Haven, Connecticut, 1985), pp. 235, 236.} It is said that Fatima's tears in Paradise for her martyred son the Imam Hussayn transformed the celestial abode of bliss into a house of mourning.\footnote{\emph{Shi`ism: Doctrines, Thought and Spirituality}, edited by Seyyed Hosein Nasr, Hamid Dabashi and Seyyed Vali Reza Nasr, Article: ``Redemptive Suffering in Islam: A Study of the Devotional Aspects of Ashura in Twelver Shi'ism'' (Albany, New York, 1988), p. 315.} The following poem was composed by Sa`id Ibn al-Nili (died 1169) for a \emph{ta`ziyeh majalis} (commemorative service for the martyrdom of Imam Hussayn). The words in the poem are supposedly spoken by Fatima Zahra:

\begin{quotex}
\emph{How great is my grief for you, o my child, you are who are the one lost of friends and family.}

\emph{Again I say how great is my sorrow, o my child, for after you I shall desert sleep and even sleeplessness.}
\emph{Woe is me, who took care of his shrouding, who beheld his face, throat and eyes.}

\emph{Woe, woe is me, who did wash him and walk behind his bier.}

\emph{Woe, woe is me, who did pray over him and lay him in his grave.}
\flright{\textsc{Ayoub}, op. cit., p. 179.}
\end{quotex}


Another source says:

\begin{quotex}
... for truly Fatima continues to weep for him (Imam Hussayn), sobbing so loudly that hell would utter such a loud cry, which, had its keepers (the angels) not been ready for it, ... its smoke and fire would have escaped and burned all that is on the face of the earth. Thus they contain Hell as long as Fatima continues to weep. ... for Hell would not calm down until her loud weeping had quieted.\footnote{Ibid., pp. 144-145.}
\end{quotex}


In \emph{Jannat al-Khulud}, Mullah Muhammad Reza Imami says that the Shi`a jinn are the ``holy jinn.'' In Iran, the pure Persian word \emph{peri} or \emph{pari} (plural: \emph{paryan}), often translated as ``fairy,'' is often used to refer to ``good jinn.'' The Arab \emph{jinn}, in reality, are radically different from the Persian \emph{pariyan.} In his \emph{Rawdat al-Shuhada}, Kashifi says that the jinn and the houris consoled and comforted Fatima and wept with her when her son the Imam Hussayn was martyred at Karbala. No doubt these were Shi'a jinn, i.e., Shi'a \emph{paryan.}

Compare the above to any of many Catholic and Eastern Orthodox songs, poems, and prayers relating to ``The Seven Sorrows of the Virgin Mary.'' Below are two examples, translations of traditional Irish-Gaelic songs dealing with the sorrows of the Virgin Mary. They are one of the many things which prompted even the irreverent non-Catholic George Bernard Shaw to exclaim: ``Holy and beautiful is the soul of Catholic Ireland.''\footnote{Excerpt from George Bernard Shaw, ``John Bull's Other Island'' in \emph{An Anthology of Irish Literature}, edited by David B. Greene (New York, 1954), p. 419.} The songs evoke both the Spanish ``saeta'' and the Urdu \emph{marsiya}, of which we shall speak in the next chapter.

\begin{quotex}
\emph{Seacht n-Dolais na Maighdine}

(The Seven Sorrows of The Virgin)

\emph{The first sorrow that was on the Virgin and she looking on her child,}

\emph{When He was born in the stable without clothes under Him nor about Him.}

(Chorus): \emph{Och ochon, Jesus, You are my child.}

\emph{Och ochon, Jesus, You are the bright king of the Heavens.}

\emph{The second sorrow that was on the Virgin and she looking on her child,}

\emph{When she got news from Egypt that their child would be taken from them.}

(Chorus)

\emph{The third sorrow that was on the Virgin and she looking on her child,}

\emph{When He was stripped of His garments and fury was put on the fair day.}

(Chorus)

\emph{The fourth sorrow that was on the Virgin and she looking on her child,}

\emph{When the crown (of thorns) was pressed on Him with spite until the blood came in a stream.}

(Chorus)

\emph{The fifth sorrow that was on the Virgin and she looking on her child,}

\emph{When He was put on the cross of torment and sharp nails binding Him.}

(Chorus)

\emph{The sixth sorrow that was on the Virgin and she looking on her child,}

\emph{When He was taken down from the cross of torment to her own bosom and He dead.}

(Chorus)

\emph{The seventh sorrow that was on the Virgin and she looking on her child,}

\emph{When He was put into the earth cold and lifeless and He dead.}

(Chorus)\footnote{Song ``\emph{Seacht n-Dolais na Maighdine}'' (\emph{The Seven Sorrows of the Virgin Mary}), LP disk Sorcha ni Ghuairim Sings Traditional Irish (Gaelic) Songs, Folkways Records Album No. FW 861. Traditional Gaelic lyrics (words of song) transcribed by Francis P. O'Connell, and also translated by him into English. New York, 1957.}

\emph{Caoineadh Mhuire} (Mary's Keen)

\emph{Peter, Apostle, have you seen my love so bright?}

\emph{M'ochon agus m'ochon o!}

\emph{I saw him with his enemies -- a harrowing sight!}

\emph{M'ochon agus m'ochon o!}

\emph{Who is that just man upon the Passion Tree?}

\emph{It is your Son, dear Mother, know you not me?}

\emph{M'ochon agus m'ochon o!}

\emph{Is that the wee babe I bore nine months in my womb}

\emph{M'ochon agus m'ochon o!}

\emph{That was born in a stable when no house would give us room?}

\emph{M'ochon agus m'ochon o!}

\emph{Mother, be quiet, let not your heart be torn}

\emph{M'ochon agus m'ochon o!}

\emph{My keening (Irish) women, mother, are yet to be born!}

\emph{M'ochon agus m'ochon o!}\footnote{\emph{``Caoineadh na Maighdine'': The Virgin's Lament}, compact disk, Noirin ni Riain and the monks of Glenstal Abbey, translations from Irish Gaelic by Gabriel Rosenstock, Glenstal Abbey, County Limerick, Ireland and Sounds True, Boulder, Colorado, 1996.
}
\end{quotex}

In Catholic and Eastern Orthodox iconography, the ``Seven Sorrows of the Virgin'' are symbolized by an image or icon of the Virgin Mary showing her heart pierced with seven swords.

A tradition related by many sources, including the highly respected historian al-Tabari, reports the following:

\begin{quotex}
Khadija, wife of the Prophet Muhammad and mother of Fatima Zahra, reported that when she was about to give birth to Fatima Zahra, the Quraishi women of Mecca, who were still pagans, refused to help her because she was the wife of Muhammad. However, the Virgin Mary, mother of Jesus Christ, called ``the prophet Isa'' by Muslims, appeared as an incredibly beautiful woman whose head was surrounded by a halo, identified herself, and proceeded to act as midwife for the delivery of Fatima Zahra.\footnote{{Abu Muhammad Ordoni, op. cit., pp. 44-45; Nasser Makarem Chirazi, op. cit., pp. 17-18; Yousuf N. Lalljee, Know Your Islam (Qum, Iran, undated), pp. 94-95.}
\end{quotex}


Below are two examples of tributes to Fatima Zahra in Shi'a literature:

\begin{quotationx}
Peace be with you, O you who were afflicted with trials by the One who created you. When He tested you, He found you to be patient under affliction... .

Peace be with you, O mistress of the women of the worlds. Peace be with you, O mother of the vindicators of humankind in argument. Peace be with you, O you who were wronged, you who were deprived of that to which you were entitled by right... .

God's blessings on the immaculate virgin, the truthful, the sinless, the pious, the unstained, the one who is pleasing to God and acceptable, the guiltless, the rightly guided, the one who was wronged, the one who was unjustly overpowered and dispossessed by force of that to which she was entitled, the one who was kept from her lawful inheritance, she whose ribs were broken, whose husband was wronged, whose son was slain, Fatima, daughter of your Prophet (Muhammad), O God, flesh of his flesh, innermost heart of his heart... mistress of women, proclaimer of God's friends, ally of piety and asceticism, apple of Paradise and Eternity... .You, O God, drew forth from her the light of the Imams.\footnote{\textsc{Abbas Qummi}, \emph{Mafatih al-Jinan} (Tehran, 1964), pp. 100-101, 539-540, 574. Cited by David Pinault in \emph{Horse of Karbala} (New York, 2001), p. 63.}

He showed them a Being, adorned with a myriad of glittering lights of various colors, who sat on a throne, a crown on her head, rings in her ears, a drawn sword by her side. The radiance streaming forth from her illumined the whole garden. When the first humans asked, ``Who is this?'', the following answer was given to them: ``This is the form of Fatima Zahra, as she appears in Paradise. Her crown is Muhammad, her earrings are Hassan and Hussayn, and her sword is `Ali (Ibn Abi Talib).\flright{\textsc{Louis Massignon}, Eranos Jahrbuch, \emph{Die Urspringe und die Bedeutung des Gnostizismus im Islam} (Zurich, 1938), pp. 64-65. Cited in \emph{Horse of Karbala} by David Pinault, op. cit., p. 65.}
\end{quotationx}

In \emph{Al-Kafi}, his monumental collection of the sayings of the Shi'a Imams, al-Kulayni cites the Imam Hussayn as commenting thusly on the ``Light Verse'' of the Qur'an:

\begin{quotex}
Abu `Abdullah (Imam Hussayn) said, concerning the words of Allah the Sublime: ``Allah is the light of the heavens and the earth. The likeness of His light is as a niche (Fatima Zahra), wherein is a lamp (al-Hassan), the lamp in a glass (al-Hussayn), the glass as it were a glittering star (Fatima Zahra, glittering star among the women of the world), kindled from a blessed tree (Ibrahim), an olive that is neither of the East nor of the West (neither of Judaism nor of Christianity), whose oil would shine (knowledge would burst out of it), even if no fire touched it. Light upon Light (Imam after Imam from the tree); Allah guides to His Light whom He wills; Allah strikes similitudes for men.\flright{\textsc{Al-Kulayni}, \emph{Al-Kafi}: volume one: \emph{Al-Usul}, part two: ``The Book of Divine Proof,'' I (Tehran, 1981), pp. 82-83.}
\end{quotex}


Commenting on the above, David Pinault says:

\begin{quotationx}
In accordance with the Shi`a tradition of viewing the Imams as the believer's means of access to God, al-Kulayni here takes the Qur'anic vocabulary of radiance and applies it to Hassan and Hussayn and their descendants. In this exegesis, the lamp-niche is allegorized as Fatima, within whom repose her sons Hassan and Hussayn, ``the lamp'' and ``the glass''. In this womb metaphor she is described as the birthplace and source of the light of the Imams.

Al-Kulayni takes this exegesis further, describing Fatima as a celestial being, foremost of the women of this lower world, linked in a kind of mystical genealogy with her spiritual forefather Abraham: starfire kindled from olivewood. Fatima the Radiant (\emph{zahra}), conveyer of illumination to her future offspring, unites celestial hierarchies, light upon light, with their earth-origins from the Abrahamic past.
\flright{\textsc{Pinault}, \emph{Horse of Karbala}, p. 65.}
\end{quotationx}

The Shi`a scholar S. V. Mir Ahmed `Ali gives a summary of the role of Fatima Zahra in Shi`ism:

\begin{quotationx}
The Holy Prophet Muhammad had two sons and a daughter, but sons he did not need, for his apostleship had to conclude with his ministry, and if any son of the Holy Prophet had survived, it would have given the chance to the people to hail the son of the Prophet also to be another prophet of God, whereas there was no prophet to come after him. It may be said that this was the reason of the male issues of the Holy Apostle leaving this world in their very infancy. But the Holy Prophet needed a daughter of his own purity of spirit and body to reflect the divine light of guidance in her ideal character and to present the authentic model of the correct Islamic womanhood to the world. Hence Fatima, the Lady of Light, was born to the Holy Prophet, who, for her godly qualities, is known in the Islamic World by the following epithets: (1) az-Zahra -- the Shining (or the Radiant), (2) al-Batul -- the Liberator of Sinners, (3) al-Azra -- the Clean, the Pure, (4) Sayyedat an-Nisa -- the Chief of Women, (5) Afzal an-Nisa -- the Most Virtuous of Women, (6) Khair an-Nisa -- the Best of Women, (7) Mariam al-Kubra -- the Supreme Mary. Mary, the mother of Jesus, was the mother of only one heavenly guide to the children of the House of Israel, whereas Lady Fatima besides being the daughter of the Holy Apostle of God, was the mother of the Eleven Divinely Commissioned Guides, (i.e., the Holy Imams)... . (8) al-Mubaraka -- the Blessed One of God, (9) as-Sadiqa -- the Truthful, (10) Al-Muhaddisa -- the one who talked to another from the womb of her mother even prior to her birth. (If God could make Jesus talk from the cradle when He was born, there should not be any wonder if the Almighty caused another one blessed by Him to act in a similar manner.) In the Qur'an it states that Jesus spoke from the womb.

There are several other epithets of this great heavenly being who is the Lady of Light. Naturally Hussayn, the King of Martyrs, could not have been born of any ordinary woman other than the one like Lady Fatima with the divine attributes she was exclusively blessed with. It has already been said that one of the unique distinctions of the family of the Holy Prophet Muhammad is that from the Holy Prophet down to the Eleventh Imam, including Lady Fatima, all the Holy Thirteen were martyrs.
\flright{\textsc{S. V. Mir Ahmed 'Ali}, \emph{Husain The Savior of Islam} (Elmhurst, New York, undated), pp. 40-41.}
\end{quotationx}

In the third paragraph of the \emph{Dawazdeh Imam} (Doxology of the Twelve Imams) by the twelfth century Persian Shi`a philosopher Nasir al-Din Tusi, he says:

\begin{quotex}
O my God! Honor and greeting, abundance and blessings be upon the glorious lady (Fatima), the beautiful, the most pure, the oppressed, the generous, the noble, who suffered so many afflictions in the course of her brief life, the queen of women, she of the great black eyes, the mother of the holy Imams, the daughter of the best of the prophets, the immaculate virgin, the most pious. Honor and salvation be upo n you and your descendants, O Fatima the radiant, O daughter of Muhammad the Messenger of God, O witness of God before His creatures. O our Lady and our Sovereign, intercede for us before God.
\flright{\textsc{Henry Corbin}, \emph{En Islam Iranien}, Tome I (Paris, 1971), p. 71.}
\end{quotex}


Here we are reminded of the Latin prayer ``\emph{Ave Maria}'' or ``Hail Mary'':

\begin{quotex}
\emph{Ave Maria, gratia plena, benedicta tu in mulieribus, et benedictus 
fructus ventre tui Jesus. Sancta Maria, Mater Dei, ora pro nobis peccatoribus, nunc et hora mortis nostrae. Amen.}

Hail Mary, full of grace, the Lord is with thee. Blessed art thou among women and blessed is the fruit of thy womb Jesus. Holy Mary, Mother of God, pray for us sinners now and at the hour of our death. Amen.
\end{quotex}


Below is a treatise on Fatima Zahra by \textbf{Henry Corbin}. As Henry Corbin makes clear, the reader is advised that what is said below is typical of the Shaikhi School of Shi`a Islam, whose views are not necessarily applicable to nor representative of Shi`ism as a whole. The reader will also note the influence on the Shaikhi School of the great Hispano-Muslim Sufi Ibn `Arabi al-Mursi and of the founder of the Illuminationist School Suhrawardi. We have noted earlier the influence of both Ibn `Arabi al-Mursi and Suhrawardi on St. John of the Cross as well. Henry Corbin's treatise proceeds as follows:

\begin{quotationx}
Perhaps we can appreciate today, even more than in the last century, philosophies that did not confuse the Imaginary, or rather the Reality corresponding to imaginative perception, with the unreal. Between a universe constituted by a pure physics and a subjectivity which inflicts isolation on itself, we foresee the need of an intermediate world to join the one with the other, something in the nature of a spiritual realm of subtle bodies. Such an intermediate world was ceaselessly meditated, particularly in Islamic Iran, by the masters of Sufism, by the adepts of the Suhrawardian philosophy of light, and by the adepts of Shaikhis. This intermediate world is no longer only the center of \emph{the world,} like Eran-Vej, but the center of \emph{the worlds.}

The world of the imaginable, of imaginative Reality, the world of archetypes-images, is established as mediator between the world of the pure, intelligible essences and the sensory universe. This world is the \emph{eighth} keshvar, the eighth climate: the ``earth of the emerald cities,'' the mystical Earth of Hurqalya... .

This is not all. When we again find Suhrawardi using the name Isfandarmuz, the Angel of the Earth and the Sophia of Mazdaism, we have no difficulty in recognizing her features, since even the characteristic name of her function has been carried over from the Mazdean liturgy into the Islamic, Neoplatonic context of Suhrawardi. But it may happen that her name is no longer pronounced, that a Figure with an entirely different name appears in an entirely different context, and that nevertheless we can still identify the same features, the same \emph{Gestalt}... .It is the feminine Archangel of a \emph{upper celestial} Earth, assuming the rank and privilege of the Divine Sophia (or \emph{Daena)}, that it is suggested that we may perceive, on the level of the world of the \emph{lahut,} the eternal reality of the dazzling Fatima, daughter of the Prophet, as she is meditated in Shi`ite gnosis, or, more exactly, in that of the Shaikhi School...

The Shaikhi School, which flourished at the end of the 18th century under the stimulus of the lofty and strong spiritual personality of Shaikh Ahmad Ahsai (d. 1826), marked an extraordinary revival of primitive Shi`ite gnosis. Its literature is enormous, for the most part still in manuscript. Here we cannot even outline all the doctrines, but in the course of the following pages we shall see how and why the theme of Hurqalya is one of its essential themes. In it, the meaning of Imamology has been closely examined in great depth. The twelve Imams who assumed the initiatic function subsequent to the prophet, his person, and the person of his daughter Fatima, from whom the line of the Imams originated, form the pleroma of the ``Fourteen Very-Pure,'' who are understood and meditated upon not only as regards the ephemeral appearance on earth of their respective persons, but in the reality of their precosmic eternal entities. Their persons are essentially theophanic; they are the Names and the divine Attributes, that which alone can be known of the divinity; they are the organs of the divinity; they are its ``operant operations''. From a structural point of view, in Shi`ite theology, Imamology plays the same role as Christology in Christain theology. That is why whoever has known only Sunni islam, is confronted in Iran by something unexpected, and becomes involved in a dialogue the richness and consequences of which are unforeseeable.

Thus, the twelve Imams, in their theophanic persons, together with the Prophet and the resplendent Fatima, form the pleroma of the ``Fourteen Very-Pure''; when meditated in their substance and their preeternal person, they assume a mode of being and a position analogous to the \emph{Aeons} of the pleroma in Valentinian gnosis. As regards the subject of our concern here, namely, the theme of the celestial Earth, the position and role of Fatima in the pleroma now take on a predominant significance. In the aforementioned schema of Suhrawardian ``Oriental Theosophy'', we were shown how our Earth and its feminine Angel, Isfandarmuz, ranked in the world of archetypes, the world of the Soul or \emph{Malakut.} Thus, we had a threefold universe: the earthly human world, which is the object of sensory perception; the world of the Soul or \emph{Malakut,} which is, properly speaking, the world of imaginative perception; and the world of pure Cherubinic Intelligences, the \emph{Jabarut}, which is the object of intelligible knowledge.

In the Shi`ite theosophy of Shaikhism, another universe (as in Ibn `Arabi al-Mursi) is superimposed on the above three universes: the universe of the \emph{lahut,} the sphere of the deity. But the characteristic of Shi`ism and Shaikhism is to conceive this \emph{lahut} explicitly as constituting the pleroma of the ``Fourteen Very-Pure''... .

In order to understand the structure of the pleroma of Shi`ite theosophy and the role played in it by Fatima, one must be guided by the basic idea... that all the universes symbolize with one another. Here again we meet the Heavens and an Earth, but these are not the Heavens and the Earth of our world, nor those of the \emph{Malakut,} nor those of the \emph{Jabarut,} but the Heavens and Earth of that \emph{hypercosmos} which is the sphere of the Deity, the \emph{lahut.} The rhythm that determines its architectonic structure is then developed in the dimension of terrestrial time. To discover in this historic dimension itself a structure which makes it possible to see the succession as homologous to the structure of the pleroma -- this will be essentially the esoteric hermeneutic, the \emph{ta'wil} -- it will be a discovery of the true and hidden meaning, the spiritual history that becomes visible through the recital of external events. It will mean to ``see things in Hurqalya''.

This question is, indeed, the key to the mystery of the primordial theophany, the revelation of the Divine Being who can only be revealed to Himself in \emph{another} self, but is unable to recognize Himself as \emph{other} or to recognize that other as Himself, except in that he \emph{Himself} is the other's God. The fact that the beings of the supreme pleroma appeared in an order of ontological precedence corresponding to the order in which they answered the primordial interrogation is a way, for the imaginative perception, of deciphering the structure of the pleroma as the place of the primordial theophany. Just as the visible Heavens are created by the contemplative acts of cherubinic Intelligences emanating one from another, so the ``heavens of the pleroma'', in the sphere of the \emph{lahut,} are brought about by theophanic acts.

These theophonic acts coincide with the progressive differentiation of the drops of the primordial ocean of being, that is, of being given its imperative by the creative \emph{Esto.} The \emph{vis formative,} immanent in each drop, enables it to give the answer that concludes the divine preeternal pact. Since the order of ontological succession of the answers determines the structure of the pleroma of the \emph{lahut,} the result is that the hierarchy of the Fourteen Supreme Spiritual Entities will have its epiphany on earth, at the time of the cycle of Muhammadan prophecy, in the succession of the persons who typify it, the ``Fourteen Very-Pure'': the Prophet Muhammad, Fatima, his daughter, and the twelve Imams.

The first of the spiritual entities to answer is the first of the beings, the ``inchoate being,'' he who will have his sensory manifestation on earth in the person of the Prophet Muhammad. This is why he is the supreme Heaven of the Pleroma, and the one whose homologue in the astronomical Heavens is the Sphere of Spheres, the Throne (\emph{'arsh}), or Empyrean. After him the second of the eternal spiritual entities to answer is the one who will be manifested on earth in the person of Hazrat Amir (i.e., the First Imam, `Ali Ibn Abi Talib, a cousin of the Prophet and the husband of Fatima); his homologue in the astronomical heavens is the Eighth Heaven, the Heaven containing the ``fortresses'' or constellations of the Zodiac, the Heaven of the Fixed Stars (\emph{Kursi}), the firmament... .

Finally there comes the response of Hazrat Fatima (Zahra) to complete the pleroma of the \emph{lahut} and give it both its plenitude and its foundation. Thus, she is the \emph{Earth} of the supreme pleroma, and this is why it can be said that on this ontological plane she is more than the Celestial Earth, she is the \emph{Supracelestial Earth.} In other words the Heavens and the Earth of the pleroma of the \emph{lahut} are related to the Heavens and the Earth of Hurqalya... in the same way as the Heavens and the Earth of Hurqalya are related to the Heavens and the Earth of the sensory world. Or again, the pleromatic person of Fatima is to the Celestial Earth of Hurqalya as Spenta Armaiti is the the Mazdean Earth haloed by the light of the \emph{Xvarnah.}

From this height, we reach a perspective in which the \emph{Sophiology} of Shaikhism will be developed. On this earth, Fatima, the daughter of the Prophet, was the wife of `Ali Ibn Abi Talib, himself the Prophet's cousin. Their exemplary union is the manifestation of the eternal syzygy originating in the eternity of the pleroma of the \emph{lahut.} The First Imam and Fatima are related to each other in the same reciprocal way as the first two hypostases, \emph{`Aql} and \emph{Nafs,} Intelligence and Soul, or in terms more familiar to us (because they go back to Philo Judaeus of Alexandria): \emph{Logos} and \emph{Sophia.}

The couple \emph{'Ali-Fatima} is the exemplification, the epiphany on earth, of the eternal couple \emph{Logos-Sophia.} Hence, we can foresee the implications of their respective persons. The Logos \emph{(`aql),} in Shaikhi doctrine is the hidden substance of every being and of every thing; it is the suprasensory calling for visible Form in order to be manifested. It is like the wood in which the form of the statue will appear. Better still, it is like the archetypal body, the inner astral mass of the sun, invisible to human perception, in relation to the visible Form, which is its \emph{aura,} brilliance and splendor. The \emph{maqam} (state, rank, degree, plane, also the pitch of a note in music) -- the \emph{maqam} of Fatima corresponds exactly to this visible form of the sun, without which there would be neither radiance nor heat. And this is why Fatima has been called by a solar name: \emph{Fatima al-Zahra,} the brilliant, resplendent Fatima. The totality of the universe consists of this light of Fatima, the splendor of each sun illuminating every conceivable universe.

So one could also speak here of a cosmic Sophianity, having its source in the eternal person of Fatima-Sophia. As such, she assumes a threefold rank, a threefold dignity and function. For she is the manifested Form, i.e., the very soul \emph{(nafs, anima)} of the Imams; she is the threshold \emph{(bab)} through which the Imams effuse the gift of their light. Just as the light of the sun is effused by the form of the sun -- which is its brilliant splendor -- not by the invisible substance of its ``archetype-body''. Thus, in the second place, she is all thinkable reality, the pleroma of meanings \emph{(ma'ani)} of all the universes, because nothing of what \emph{is} can be without qualification. Qualification and meaning are on the same level of being as the Soul, for it is Soul-Sophia that confers qualification and meaning. That is why the whole universe of the soul and the secret of the meanings given by the Soul is the very universe and secret of Hazrat Fatima. She is Sophia, which is to say divine wisdom and power, embracing all the universes. That, lastly, is why her eternal Person, which is the secret of the world of the Soul, is also its manifestation \emph{(bayan),} without which the creative Principle of the world would remain unknown and unknowable, forever hidden... .

Fatima-Sophia is in fact the Soul: the Soul of creation, the Soul of each creature, i.e., the constitutive part of the human being that appears essentially to the imaginative consciousness in the form of a feminine being, \emph{anima.} She is the eternally feminine in man, and that is why she is the archetype of the heavenly Earth; she is both paradise and initiation into it, for it is she who manifests the divine names and attributes revealed in the theophanic persons of the Imams, i.e., in the Heavens of the Pleroma of the \emph{lahut}... .

Her functions \emph{symbolize} with each other, from one universe to the other: in the pleroma of the \emph{lahut,} as the supracelestial Earth which is its foundation; on the terrestrial Earth, as the daughter and Soul of the Prophet and as the one from whom issue those who in their turn are the Soul of the Prophet, the lineage of the Twelve Imams. She is \emph{the} theophany and she is the Initiation; she is the \emph{majma alnurayn,} the confluence of two lights, the light of Prophecy and the light of Initiation. Through her, creation, from the beginning, is Sophianic in nature, and through her the Imams are invested with the Sophianity that they transmit to their adepts, because she is its \emph{soul.} From this pleromic height we can distinguish the fundamental sound emerging from the depths: namely, that which Mazdean Sophiology formulated in the idea of \emph{spendar matikih,} the Sophianity with which Spenta Armaiti, the feminine Angel of the Earth invested the faithful believer.
\flright{\textsc{Henry Corbin}, \emph{Spiritual Body and Celestial Earth}, translated by Nancy Pearson (Princeton, New Jersey, 1977), pp. 50-73.}
\end{quotationx}

\flrightit{Posted on 2016-09-06 by Michael McClain}

\begin{center}* * *\end{center}

\begin{footnotesize}
\begin{sffamily}
\texttt{Matt on 2016-09-06 at 23:52 said:}

The way theology is done in Shia Islam also has an affinity with the Catholic Church. Philosophical theology has an important place: a notable example of this is seen with Ijtihad (independant reasoning), which is used a lot in jurispudence among marja and lower-level Shia clerics.


\hfill

\end{sffamily}
\end{footnotesize}

